\begin{center}
    \textbf{\LARGE{Devoir maison : vérifier qu'un tableau est trié}}\\
\end{center}


On s'intéresse à la fonction \textsc{verifie} qui prend en paramètre un tableau de valeurs numériques non vide et un paramètre ordre permettant de choisir l'ordre dans lequel le tableau doit être trié. La fonction renvoie True si ce tableau est trié dans l’ordre donné, False sinon.

\begin{lstlisting}[language=Python]
def verifie(tab, ordre="croissant"):

    # Tri du plus petit au plus grand
    if ordre == "croissant" :
        for i in range(len(tab)-1):
            if tab[i] > tab[i+1]:
                return False
            return True
    
    # Tri du plus grand au plus petit    
    else : 
            for i in range(len(tab)-1):
                if _ _ _ _ _ _ _ _ _ _ :
                    return _ _ _ _ 
                return _ _ _ _ 
\end{lstlisting}

\begin{Exercise}[title={Vérification tri}]
\Question Combien y a-t-il de \textbf{conditions} dans le code ci-dessus ?\\
\emptybox{1cm}
\Question Quel \textbf{type de parcours} est utilisé à la ligne 3 de la fonction ?\\
\emptybox{1cm}
\Question Complétez les trous dans le programme.\\
\Question Compléter le tableau d'avancement ci-dessous pour un appel de la fonction suivant : \textsc{verifie([0, 5, 8, 8, 9])}.\\
\begin{center}
\large
\begin{tabular}{c|c|c|c|c|}
\textbf{Tour} & \textbf{i} & \textbf{tab[i]} & \textbf{tab[i+1]} & \textbf{tab[i] > tab[i+1]}  \\
\hline
0 &  &  &  & \\
1 &  &  &  & \\
2 &  &  &  & \\
3 &  &  &  & \\
4 &  &  &  & \\
\end{tabular}
\end{center}

\Question Quelles sont les valeurs de retours de la fonction verifie pour les appels suivants :
\begin{itemize}
 \item verifie([-1, 4, 8, 8, 88])
 \item verifie([5, 3, 1], ordre="decroissant")
 \item verifie([5, 3, 1])
 \item verifie([7, 8, 9], ordre="decroissant")
\end{itemize}


\Question Question BONUS : Proposez une implémentation plus concise de la fonction verifie.
\end{Exercise}